% BEGIN VERIFIED DISPLAY
The only unknowns are the forty separately named positive-integer variables
\[
\begin{gathered}J,F,\alpha,f,g,T,X,Y,\xi,\eta,\zeta,\upsilon,\gamma,\vartheta,\psi,\varphi,\epsilon,\\
\ell,W,s,x,m,n,z,d,V,\delta,\\
\kappa,\lambda,\tau,\theta,\nu,Q,\\
N,O,P,S,A,i,\Gamma\in\mathbb Z_{>0}.\end{gathered}
\]
All other symbols below are polynomial abbreviations:
\[
\begin{gathered}
b=\beta J,\quad q=1+b,\quad C=q-F-\alpha,\quad U=Y\kappa,\\
c=\eta+O,\quad a=\lambda+Y,\\
E=\lambda,\quad k=\eta+\zeta,\\
G=i,\\
H=3+4a,\quad\Delta=\delta,\quad L=P,\\
M=sb,\\
R=G(\kappa-1)+J(\mu+\omega q),\\
I=\theta,\quad Z=\psi z,\quad D=\varphi d,\\
\mathcal{P}=z^2-(1+\tau),\\
\mathcal{D}=d^2-(1+N),\\
\mathcal{Q}=\nu-Q[\nu-\vartheta]-1,\\
\mathcal{R}=I^2-Q,\quad \mathcal{E}=\vartheta-V^2,\\
\mathcal{B}=\psi-M,\quad \mathcal{C}=\varphi-U,\quad \mathcal{G}=\epsilon-R,\\
\mathcal{L}=\ell-EO,\quad \mathcal{W}=W-[C(X+Kq)-Fq],\\
\mathcal{X}=s-q\kappa,\quad \mathcal{Y}=x-Eq,\\
\mathcal{Z}=m-f^2,\quad \mathcal{N}=n-c^2,\\
\mathcal{I}=Z-[\psi(\ell+k+\xi-\Gamma)+Yb-Ws],\\
\mathcal{J}=D-[\varphi L+H(S+x+\epsilon+1-k)],\\
\mathcal{K}=V-[m(A-\epsilon)-c],\quad\mathcal{O}=\delta-H-a^2,\\
\mathfrak{J}=\kappa-q^2,\quad\mathcal{F}=\lambda-XY,\\
\mathcal{S}=\tau-\ell(\ell+k),\\
\mathcal{V}=\theta-ng,\\
\mathfrak{A}=\nu-\upsilon^2,\quad\mathfrak{B}=Q-\Delta(m-1),\\
\mathfrak{C}=N-\delta n,\quad\mathfrak{D}=G-[\alpha+(q-F)b],\\
\mathfrak{E}=O-Yk,\quad\mathfrak{F}=\Gamma-kW,\\
\mathfrak{G}=P-X-ac,\quad\mathfrak{H}=S-\varphi\gamma,\\
\mathfrak{I}=A-Tc.
\end{gathered}
\]
The single Diophantine equation is
\[
\begin{split}
\mathcal{P}^2+\mathcal{D}^2+\mathcal{Q}^2+\mathcal{R}^2+\mathcal{E}^2+\mathcal{B}^2+\mathcal{C}^2+\mathcal{G}^2 &\\
+\mathcal{L}^2+\mathcal{W}^2+\mathcal{X}^2+\mathcal{Y}^2+\mathcal{Z}^2+\mathcal{N}^2 &\\
+\mathcal{I}^2+\mathcal{J}^2+\mathcal{K}^2+\mathcal{O}^2+\mathfrak{J}^2+\mathcal{F}^2 &\\
+\mathcal{S}^2+\mathcal{V}^2+\mathfrak{A}^2+\mathfrak{B}^2 &\\
+\mathfrak{C}^2+\mathfrak{D}^2+\mathfrak{E}^2+\mathfrak{F}^2+\mathfrak{G}^2+\mathfrak{H}^2+\mathfrak{I}^2 &= 0.
\end{split}
\]
% END VERIFIED DISPLAY
